Let $D$ be the spring constant. The total energy stays the same: at a turning point the body is at rest, and all of the energy is potential:
\begin{align}
 E &= \tfrac12\,D\,A^2
\end{align}
At a displacement $y$ the potential energy is $E_\mathrm{pot} = \tfrac12\,D\,y^2$ (it grows with $y^2$, not with $y$), and the kinetic energy is the rest, $E_\mathrm{kin} = E - E_\mathrm{pot}$.

\begin{abcliste}
\abc
\begin{align}
 \frac{E_\mathrm{pot}}{E} &= \frac{\tfrac12\,D\,y^2}{\tfrac12\,D\,A^2} = \potF = \left(\frac{3}{5}\right)^2 = \frac{9}{25} = \pot \\
 \frac{E_\mathrm{kin}}{E} &= \kinF = 1-\frac{9}{25} = \frac{16}{25} = \kin \approx \result{\kinP{0}{2}}
\end{align}
At $\frac35$ of the amplitude the body still has almost two thirds of its energy as kinetic energy.

\abc
Equal kinetic and potential energy means each is half of $E$:
\begin{align}
 \tfrac12\,D\,y^2 &= \tfrac12\cdot\tfrac12\,D\,A^2 \quad\Rightarrow\quad y = \frac{A}{\sqrt2}
\end{align}
\begin{align}
 \frac{y}{A} &= \yeF = \ye \approx \result{\yeP{0}{2}}
\end{align}
So $y \approx 0.71\,A$, not $\frac12 A$: at half the amplitude the potential energy is only $\left(\frac12\right)^2 = \frac14$ of the total.
\end{abcliste}
